\DocumentMetadata{
  pdfversion=2.0,
  pdfstandard=ua-2,
  lang=en-US,
  tagging=on,
  tagging-setup={math/setup=mathml-SE,tabsorder=structure}
}
\documentclass[11pt,letterpaper]{article}
\usepackage[margin=0.68in,includefoot]{geometry}
\usepackage{newcomputermodern}
\usepackage{amsmath,amscd,mathtools}
\usepackage{tikz,adjustbox}
\providecommand{\square}{\mdlgwhtsquare}
\usepackage[hidelinks]{hyperref}
\hypersetup{
  pdftitle={Algebra First-Year Exam, May 2020, Part 2},
  pdfauthor={University of Florida Department of Mathematics},
  pdfsubject={Graduate mathematics examination},
  pdfdisplaydoctitle=true,
  bookmarksopen=true
}
\setcounter{secnumdepth}{0}
\setlength{\parindent}{0pt}
\setlength{\parskip}{0.28em}
\setlength{\emergencystretch}{3em}
\setlength{\leftmargini}{2.15em}
\setlength{\leftmarginii}{2.65em}
\setlength{\leftmarginiii}{3em}
\renewcommand{\labelenumi}{\textbf{\theenumi.}}
\pagestyle{empty}
\raggedbottom
\allowdisplaybreaks
\newenvironment{examproblems}[1]{%
  \begin{enumerate}%
  \setcounter{enumi}{#1}%
  \setlength{\topsep}{0.35em}%
  \setlength{\partopsep}{0pt}%
  \setlength{\itemsep}{0.48em}%
  \setlength{\parsep}{0.18em}%
}{\end{enumerate}}
\newenvironment{examsubparts}{%
  \begin{enumerate}%
  \setlength{\topsep}{0.18em}%
  \setlength{\partopsep}{0pt}%
  \setlength{\itemsep}{0.16em}%
  \setlength{\parsep}{0.1em}%
}{\end{enumerate}}
\newenvironment{examinstructions}{%
  \par\begingroup\itshape
}{%
  \par\endgroup\vspace{0.18em}
}
\makeatletter
\renewcommand\section{\@startsection{section}{1}{0pt}%
  {-1.25ex plus -.25ex minus -.1ex}%
  {0.55ex plus .1ex}%
  {\normalfont\large\bfseries}}
\renewcommand{\@maketitle}{%
  \newpage\null\vskip -1em%
  {\footnotesize\bfseries
    \href{https://www.ufl.edu/}{University of Florida}, \href{https://math.ufl.edu/}{Department of Mathematics}\par}%
  \vskip -0.5em%
  \tagstructbegin{tag=Title}%
    {\LARGE\bfseries\href{https://gma.math.ufl.edu/exams/algebra-first-year/}{Algebra First-Year Exam}, May 2020, Part 2\par}%
  \tagstructend%
  \vskip 0.25em%
  \tagmcbegin{artifact}%
    \hrule height 1pt%
  \tagmcend%
  \vskip 1em%
}
\makeatother
\title{Algebra First-Year Exam, May 2020, Part 2}
\author{}
\date{}
\begin{document}
\maketitle
\thispagestyle{empty}
\begin{examinstructions}
Do four out of the six problems below. If you hand in more than four problems, only the first four will be graded. You may (within reason) use results from the text if you state them fully.
\end{examinstructions}
\begin{examproblems}{0}
\item
Suppose~\(R\) is a ring with identity, and recall that a nonzero unital (left)~\(R\)-module~\(M\) is simple if its only submodules are~\(0\) and~\(M\). Suppose that~\(M\) is a simple~\(R\)-module.
\begin{examsubparts}
\item[\textnormal{(i)}]
Show that~\(M\) is cyclic, i.e. generated by a single element.
\item[\textnormal{(ii)}]
Show that the ring \(\operatorname{End}(M)\) of all~\(R\)-linear endomorphisms of~\(M\) is a division ring.
\item[\textnormal{(iii)}]
Show that if~\(R\) is commutative then \(M \simeq R/\mathfrak{m}\) for some maximal ideal \(\mathfrak{m} \subset R\), and identify \(\operatorname{End}(M)\) in this case.
\end{examsubparts}
\item
Which of the following polynomials are irreducible in the given ring? Explain briefly, and if the polynomial is not irreducible, give an irreducible factorization.
\begin{examsubparts}
\item[\textnormal{(i)}]
\(X^3+5X^2+8X+6 \in \mathbb{Q}[X]\)
\item[\textnormal{(ii)}]
\(X^5+(2+3i)X+13 \in \mathbb{Q}(i)[X]\)
\item[\textnormal{(iii)}]
\(X^3+Y^3-1 \in \mathbb{Q}[X,Y]\)
\end{examsubparts}
\item
Let~\(R\) be the ring
\[
R=\{a+b\sqrt{-13}\in\mathbb{C}\mid a,b\in\mathbb{Z}\}.
\]
\begin{examsubparts}
\item[\textnormal{(i)}]
Find the units of~\(R\) (introduce an appropriate norm).
\item[\textnormal{(ii)}]
Show that \(14\in R\) has two distinct factorizations into irreducible elements.
\item[\textnormal{(iii)}]
Show that \((7,1+\sqrt{-13})\) and \((7,1-\sqrt{-13})\) are prime ideals in~\(R\).
\item[\textnormal{(iv)}]
Show that \((7)\) is not a prime ideal in~\(R\), and write it as a product of prime ideals.
\end{examsubparts}
\item
Let~\(K\) be a field. Suppose~\(V\) is a~\(K\)-vector space of finite dimension and that \(W_1,W_2\subseteq V\) are subspaces. Show that
\[
\dim_K(W_1+W_2)=\dim_K(W_1)+\dim_K(W_2)-\dim_K(W_1\cap W_2).
\]
\item
Find representatives in rational canonical form for all conjugacy classes of order~\(5\) or~\(7\) in \(GL_4(\mathbb{F}_2)\) (first find irreducible factorizations of \(X^5-1\) and \(X^7-1\) in \(\mathbb{F}_2[X]\)).
\item
Suppose~\(p\) and~\(q\) are distinct primes in~\(\mathbb{Z}\) and let \(\alpha=\sqrt{p}+\sqrt{q}\). Show that \(\mathbb{Q}(\alpha)=\mathbb{Q}(\sqrt{p},\sqrt{q})\) and use this to compute the degree of the extension \(\mathbb{Q}(\alpha)/\mathbb{Q}\). Find also the minimal polynomial of~\(\alpha\) over~\(\mathbb{Q}\).
\end{examproblems}
\end{document}
